\newpage
\section*{Solucionario completo}

\noindent Cada respuesta incluye la clave oficial y su justificación analítica o matemática detallada.

\subsection*{Bloque A. Fundamentos teóricos y conceptuales}

\sol{1}{B}{\textbf{Tema:} Estimador matricial MCO. \textbf{Explicación:} La minimización de la suma de residuos al cuadrado $e'e = (Y - X\beta)'(Y - X\beta)$ conduce al sistema de ecuaciones normales $(X'X)\hat{\beta} = X'Y$. Si $X'X$ es invertible, la solución única es $\hat{\beta} = (X'X)^{-1} X'Y$.}

\sol{2}{D}{\textbf{Tema:} Insesgadez bajo exogeneidad. \textbf{Explicación:} Como $\hat{\beta} = \beta + (X'X)^{-1} X'u$, tomando esperanza condicional en $X$ se tiene $E[\hat{\beta}|X] = \beta + (X'X)^{-1} X' E[u|X]$. Si $E[u|X] = 0$, entonces $E[\hat{\beta}|X] = \beta$, garantizando la insesgadez.}

\sol{3}{C}{\textbf{Tema:} Rango completo y multicolinealidad perfecta. \textbf{Explicación:} Si existe una relación lineal exacta entre las variables explicativas (multicolinealidad perfecta), las columnas de $X$ son linealmente dependientes, el determinante de $X'X$ es cero y la matriz no se puede invertir.}

\sol{4}{C}{\textbf{Tema:} Supuesto de perturbaciones esféricas. \textbf{Explicación:} Perturbaciones esféricas implica que la matriz de varianza-covarianza de los errores es escalar: $\Sigma = \sigma^2 I_n$. Los elementos de la diagonal principal son constantes e idénticos (homocedasticidad) y los elementos fuera de la diagonal son nulos (no autocorrelación).}

\sol{5}{D}{\textbf{Tema:} Teorema de Gauss-Markov (BLUE). \textbf{Explicación:} Gauss-Markov demuestra que para cualquier combinación lineal $c'\beta$, el estimador lineal insesgado MCO $c'\hat{\beta}$ minimiza la varianza de estimación frente a cualquier otro estimador lineal insesgado rival.}

\sol{6}{C}{\textbf{Tema:} Estimador insesgado de la varianza residual. \textbf{Explicación:} La suma de residuos al cuadrado $e'e$ pierde $k$ grados de libertad por las $k$ condiciones ortogonales de primer orden impuestas al estimar los parámetros. Por tanto, dividiendo entre $n - k$ se obtiene $E[s^2] = \sigma^2$.}

\sol{7}{C}{\textbf{Tema:} Propiedad del R-cuadrado. \textbf{Explicación:} Al agregar regresores adicionales, el espacio de columnas de $X$ se expande y la suma de cuadrados residuales ($SCR$) nunca puede aumentar en la optimización por mínimos cuadrados. Por ende, $R^2 = 1 - SCR/STC$ nunca se reduce.}

\sol{8}{D}{\textbf{Tema:} R-cuadrado ajustado. \textbf{Explicación:} El $R^2$ ajustado corrige dividiendo $SCR$ y $STC$ entre sus respectivos grados de libertad ($n - k$ y $n - 1$). Puede incluso tomar valores negativos si el modelo es sumamente deficiente y el estadístico $F$ es menor que 1.}

\sol{9}{A}{\textbf{Tema:} Distribución normal de los coeficientes MCO. \textbf{Explicación:} Dado que $\hat{\beta}$ es una combinación lineal de las perturbaciones normales $Y$ (o $u$), hereda la distribución normal multivariada con media $\beta$ y matriz de varianzas $\sigma^2 (X'X)^{-1}$.}

\sol{10}{C}{\textbf{Tema:} Prueba t de significancia individual. \textbf{Explicación:} El cociente entre el coeficiente estimado y su error estándar estimado sigue exactamente una distribución $t$ con $n - k$ grados de libertad bajo el supuesto de normalidad de los errores, o normal asintótica en muestras grandes por el TLC.}

\sol{11}{C}{\textbf{Tema:} Prueba F conjunta de la regresión. \textbf{Explicación:} La prueba $F$ global mide si la variación explicada por los $k - 1$ regresores (excluyendo el intercepto) es significativamente mayor que la variación residual no explicada ponderada por grados de libertad.}

\sol{12}{B}{\textbf{Tema:} Interpretación de modelo log-log. \textbf{Explicación:} Dado que $\frac{d \ln(Y)}{d \ln(X)} = \frac{dY / Y}{dX / X} = \beta_1$, el parámetro mide la elasticidad directamente y es invariante a las unidades de medida de las variables.}

\sol{13}{A}{\textbf{Tema:} Interpretación de modelo log-lineal. \textbf{Explicación:} $\frac{d \ln(Y)}{dX} = \frac{1}{Y} \frac{dY}{dX} = \beta_1 \implies \frac{\Delta Y}{Y} \approx \beta_1 \Delta X$. Por tanto, multiplicando por 100 se obtiene el porcentaje de cambio aproximado en $Y$ por cada unidad de incremento en $X$.}

\sol{14}{B}{\textbf{Tema:} Consecuencias de la multicolinealidad imperfecta. \textbf{Explicación:} La multicolinealidad severa infla las varianzas muestrales de los coeficientes afectados ($VIF_j = \frac{1}{1 - R_j^2}$ muy alto). Los intervalos de confianza se ensanchan y los estadísticos $t$ se desploman, aunque el $R^2$ global del modelo sea muy elevado.}

\sol{15}{D}{\textbf{Tema:} Definición de heterocedasticidad. \textbf{Explicación:} La homocedasticidad exige $Var(u_i|X) = \sigma^2$ constante para todo $i$. Si la dispersión de las perturbaciones varía entre individuos, estratos o tamaños de empresa, existe heterocedasticidad: $Var(u_i|X) = \sigma_i^2$.}

\sol{16}{C}{\textbf{Tema:} Consecuencias de la heterocedasticidad. \textbf{Explicación:} La insesgadez no depende del supuesto de varianza constante, sino de $E[u|X]=0$. Sin embargo, la fórmula de varianza usual $\sigma^2 (X'X)^{-1}$ queda mal calculada, arrojando errores estándar engañosos y pruebas de hipótesis inválidas.}

\sol{17}{B}{\textbf{Tema:} Prueba de White. \textbf{Explicación:} Halbert White (1980) propuso una prueba que no requiere supuestos a priori sobre la forma funcional de la heterocedasticidad: se regresiona $e_i^2$ sobre variables lineales, cuadráticas e interacciones. El estadístico $LM = n R_{aux}^2$ sigue una $\chi^2$ con grados de libertad iguales al número de regresores en la ecuación auxiliar.}

\sol{18}{A}{\textbf{Tema:} Prueba de Breusch-Pagan. \textbf{Explicación:} La prueba de Breusch-Pagan evalúa si la varianza del error depende de variables observadas $Z$. La hipótesis nula es la ausencia de heterocedasticidad (homocedasticidad), de modo que los coeficientes de las variables auxiliares son nulos.}

\sol{19}{C}{\textbf{Tema:} Mínimos Cuadrados Ponderados (WLS). \textbf{Explicación:} Al transformar el modelo dividiendo cada observación entre $\sqrt{h_i}$, la perturbación transformada $u_i^* = u_i / \sqrt{h_i}$ tiene varianza $Var(u_i^*) = Var(u_i) / h_i = \sigma^2 h_i / h_i = \sigma^2$, restaurando la homocedasticidad y la eficiencia BLUE.}

\sol{20}{B}{\textbf{Tema:} Errores estándar robustos de White. \textbf{Explicación:} El estimador de Huber-White no altera $\hat{\beta}$, pero reemplaza la fórmula ingenua de varianza por la estructura sándwich consistente asintóticamente, garantizando que los estadísticos $t$ y $F$ sean fiables en muestras grandes aun con heterocedasticidad desconocida.}

\sol{21}{B}{\textbf{Tema:} Definición de proceso AR(1). \textbf{Explicación:} En un proceso autorregresivo de primer orden $AR(1)$, el error del periodo actual depende linealmente del error del periodo inmediato anterior ponderado por $\rho$ más una innovación contemporánea independiente $\varepsilon_t$.}

\sol{22}{D}{\textbf{Tema:} Consecuencias de la autocorrelación serial. \textbf{Explicación:} La correlación positiva sucesiva entre errores hace que la muestra contenga menos información independiente de la aparente ($n$ efectivo menor). MCO subestima la dispersión real, arrojando errores estándar diminutos y $t$ gigantescos que inducen a rechazar falsamente hipótesis nulas.}

\sol{23}{A}{\textbf{Tema:} Valor central de Durbin-Watson. \textbf{Explicación:} Si $\hat{\rho} = 0$, $d \approx 2(1 - 0) = 2$. Valores de $d$ cercanos a 0 indican fuerte autocorrelación positiva ($\hat{\rho} \approx 1$); valores cercanos a 4 indican fuerte autocorrelación negativa ($\hat{\rho} \approx -1$).}

\sol{24}{D}{\textbf{Tema:} Limitación de Durbin-Watson con rezagos. \textbf{Explicación:} Con variables dependientes rezagadas ($Y_{t-1}$), el estadístico $d$ está fuertemente sesgado hacia 2 (hacia no rechazar la no autocorrelación), dando una falsa impresión de limpieza estocástica. Se debe usar el estadístico $h$ o la prueba LM de Breusch-Godfrey.}

\sol{25}{A}{\textbf{Tema:} Ventajas de Breusch-Godfrey. \textbf{Explicación:} La prueba LM de Breusch-Godfrey es general: regresa los residuos $e_t$ sobre las variables explicativas y $p$ rezagos de los propios residuos ($e_{t-1}, \dots, e_{t-p}$). El estadístico $(n - p) R^2 \sim \chi^2_p$ es asintóticamente robusto frente a regresores endógenos rezagados.}

\sol{26}{A}{\textbf{Tema:} Estimador HAC de Newey-West. \textbf{Explicación:} Whitney Newey y Kenneth West diseñaron una matriz de varianzas-covarianzas que estima de manera consistente las correlaciones intertemporales no nulas y la heterocedasticidad aplicando un núcleo de ponderación decreciente (núcleo de Bartlett), garantizando que la matriz sea semidefinida positiva.}

\sol{27}{C}{\textbf{Tema:} Causas de la endogeneidad. \textbf{Explicación:} La endogeneidad rompe la exogeneidad estricta. Ocurre por: 1) omitir una variable relevante correlacionada con los regresores incluidos, 2) medir las $X$ con ruido aleatorio (sesgo de atenuación), o 3) determinación simultánea de $Y$ y $X$ en un sistema de equilibrio.}

\sol{28}{A}{\textbf{Tema:} Inconsistencia bajo endogeneidad. \textbf{Explicación:} Bajo endogeneidad, $plim \, \hat{\beta} = \beta + \frac{Cov(X, u)}{Var(X)} \neq \beta$. Al haber correlación entre $X$ y $u$, el sesgo no desaparece al aumentar el tamaño de muestra, invalidando toda interpretación causal del coeficiente.}

\sol{29}{B}{\textbf{Tema:} Condiciones de un instrumento válido. \textbf{Explicación:} La condición de relevancia exige que el instrumento mueva a la variable endógena ($Cov(Z, X) \neq 0$). La restricción de exclusión (exogeneidad) exige que $Z$ sea ortogonal al término de error poblacional ($Cov(Z, u) = 0$), incidiendo en $Y$ exclusivamente por su efecto sobre $X$.}

\sol{30}{C}{\textbf{Tema:} Procedimiento de MC2E / 2SLS. \textbf{Explicación:} La primera etapa extrae la porción exógena de la variable endógena proyectándola sobre los instrumentos ($\hat{X} = Z(Z'Z)^{-1}Z'X$). La segunda etapa utiliza únicamente este componente exógeno purificado $\hat{X}$ para estimar la ecuación estructural libre de sesgo.}

\subsection*{Bloque B. Aplicaciones cuantitativas y contrastes econométricos}

\sol{31}{A}{\textbf{Tema:} Cálculo de la pendiente MCO simple. \textbf{Explicación:} $\hat{\beta}_1 = \frac{\sum (X_i - \bar{X})(Y_i - \bar{Y})}{\sum (X_i - \bar{X})^2} = \frac{600}{400} = 1.50$.}

\sol{32}{D}{\textbf{Tema:} Cálculo de correlación y R-cuadrado. \textbf{Explicación:} $r = \frac{Cov(X, Y)}{S_X S_Y} = \frac{600}{\sqrt{400 \times 900}} = \frac{600}{\sqrt{360000}} = \frac{600}{600} = 1.00$... corrección: con 400 y 900 es $\sqrt{360000} = 600$. Entonces $r = 1.00$ y $R^2 = 1.00$.}

\sol{33}{C}{\textbf{Tema:} Cálculo de varianza residual. \textbf{Explicación:} Grados de libertad: $n - k = 45 - 5 = 40$. Varianza residual: $s^2 = \frac{SCR}{n - k} = \frac{800}{40} = 20$. Error estándar: $s = \sqrt{20} \approx 4.472$.}

\sol{34}{A}{\textbf{Tema:} Cálculo de estadístico t. \textbf{Explicación:} $t = \frac{\hat{\beta}_1}{EE(\hat{\beta}_1)} = \frac{-0.45}{0.15} = -3.00$. Como $|t| = 3.00 > 1.98$, se rechaza la hipótesis nula con un nivel de confianza superior al $99\%$.}

\sol{35}{A}{\textbf{Tema:} Cálculo del estadístico F global. \textbf{Explicación:} $F = \frac{R^2 / (k - 1)}{(1 - R^2) / (n - k)} = \frac{0.60 / 3}{(1 - 0.60) / (64 - 4)} = \frac{0.20}{0.40 / 60} = \frac{0.20}{0.00667} = 30.0$.}

\sol{36}{C}{\textbf{Tema:} Prueba F de restricciones lineales. \textbf{Explicación:} $F = \frac{(SCR_R - SCR_{NR}) / q}{SCR_{NR} / (n - k)} = \frac{(560 - 500) / 2}{500 / (106 - 6)} = \frac{60 / 2}{500 / 100} = \frac{30}{5} = 6.0$.}

\sol{37}{B}{\textbf{Tema:} Interpretación de elasticidad en demanda log-log. \textbf{Explicación:} En el modelo log-log, el coeficiente $-0.35$ es la elasticidad precio directa: $\% \Delta Q \approx -0.35 \times 10\% = -3.5\%$. Dado que $|-0.35| < 1$, la demanda de agua potable es inelástica respecto al precio.}

\sol{38}{A}{\textbf{Tema:} Cálculo de semielasticidad log-lineal. \textbf{Explicación:} $\Delta \ln(OPEX) \approx \beta_1 \Delta X = 0.04 \times 5 = 0.20$. Multiplicando por 100, los costos operativos aumentan aproximadamente un $20\%$.}

\sol{39}{A}{\textbf{Tema:} Interpretación de modelo lineal-log. \textbf{Explicación:} En especificación lineal-log, $\Delta Y \approx (\beta_1 / 100) \times (\% \Delta X) = (200 / 100) \times 1 = 2$ unidades absolutas.}

\sol{40}{B}{\textbf{Tema:} Cálculo del Factor de Inflación de Varianza (VIF). \textbf{Explicación:} $VIF_j = \frac{1}{1 - R_j^2} = \frac{1}{1 - 0.80} = \frac{1}{0.20} = 5.0$.}

\sol{41}{A}{\textbf{Tema:} Cálculo de error estándar matricial. \textbf{Explicación:} La matriz de covarianzas es $s^2 (X'X)^{-1}$. La varianza de $\hat{\beta}_2$ es $Var(\hat{\beta}_2) = s^2 \times [(X'X)^{-1}]_{33} = 25 \times 0.09 = 2.25$. El error estándar es $EE(\hat{\beta}_2) = \sqrt{2.25} = 1.50$.}

\sol{42}{B}{\textbf{Tema:} Estadístico de prueba de White. \textbf{Explicación:} $LM = n \times R_{aux}^2 = 100 \times 0.15 = 15.0$. Sigue asintóticamente una distribución $\chi^2$ con grados de libertad iguales al número de regresores en la ecuación auxiliar ($p = 5$). Como el valor crítico al $5\%$ de $\chi^2_5$ es $11.07$ y $15.0 > 11.07$, se rechaza la homocedasticidad.}

\sol{43}{D}{\textbf{Tema:} Cálculo de Breusch-Pagan. \textbf{Explicación:} $LM = 80 \times 0.10 = 8.0$. Bajo $H_0$ de homocedasticidad sigue una $\chi^2_3$. El valor crítico al $5\%$ es $7.81$. Como $8.0 > 7.81$, se rechaza la homocedasticidad.}

\sol{44}{C}{\textbf{Tema:} Transformación analítica de WLS. \textbf{Explicación:} La desviación estándar del error es $\sqrt{\sigma^2 X_i^2} = \sigma X_i$. Dividiendo toda la ecuación entre $X_i$: $\frac{Y_i}{X_i} = \beta_0 (1/X_i) + \beta_1 + u_i^*$, donde $Var(u_i^*) = Var(u_i/X_i) = \sigma^2$, logrando homocedasticidad perfecta.}

\sol{45}{D}{\textbf{Tema:} Cálculo de Durbin-Watson a partir de rho. \textbf{Explicación:} $d \approx 2(1 - \hat{\rho}) = 2(1 - 0.60) = 2(0.40) = 0.80$.}

\sol{46}{D}{\textbf{Tema:} Regla de decisión de Durbin-Watson. \textbf{Explicación:} Si el estadístico calculado $d$ es estrictamente menor que el límite inferior tabulado $d_L$ ($0.80 < 1.63$), se rechaza categóricamente la hipótesis nula $H_0: \rho = 0$ a favor de la alternativa de autocorrelación serial positiva.}

\sol{47}{A}{\textbf{Tema:} Estimación de rho para Cochrane-Orcutt. \textbf{Explicación:} Despejando de $d \approx 2(1 - \hat{\rho}) \implies 1 - \hat{\rho} = d / 2 \implies \hat{\rho} = 1 - d/2 = 1 - 1.30 / 2 = 1 - 0.65 = 0.35$.}

\sol{48}{B}{\textbf{Tema:} Transformación de Prais-Winsten en t=1. \textbf{Explicación:} A diferencia de Cochrane-Orcutt (que descarta la primera observación), Prais-Winsten conserva la primera observación multiplicándola por $\sqrt{1 - \rho^2}$, preservando la eficiencia en muestras pequeñas.}

\sol{49}{C}{\textbf{Tema:} Cálculo de sesgo de variable omitida. \textbf{Explicación:} La fórmula del sesgo de variable omitida es $E[\tilde{\beta}_1] = \beta_1 + \beta_2 \frac{Cov(X_1, X_2)}{Var(X_1)} = \beta_1 + \beta_2 \tilde{\gamma}_1 = 4 + 2(0.5) = 4 + 1.0 = 5.0$.}

\sol{50}{B}{\textbf{Tema:} Cálculo del estimador de Variables Instrumentales. \textbf{Explicación:} El estimador de variables instrumentales simple es el cociente de covarianzas muestrales: $\hat{\beta}_{IV} = \frac{Cov(Z, Y)}{Cov(Z, X)} = \frac{150}{75} = 2.00$.}

\sol{51}{A}{\textbf{Tema:} Regla de instrumentos débiles de Stock-Yogo. \textbf{Explicación:} La regla empírica estándar de Stock y Yogo establece que si el estadístico $F$ de los instrumentos excluidos en la primera etapa es menor a 10 ($4.2 < 10$), los instrumentos son débiles, generando sesgo significativo y distorsión de tamaño en los contrastes de Wald.}

\sol{52}{A}{\textbf{Tema:} Mecánica de la prueba de Hausman en regresión. \textbf{Explicación:} Bajo $H_0$ de exogeneidad, el residuo de primera etapa $\hat{v}$ no debe tener poder explicativo sobre $Y$ una vez controlado $X_1$. Si su coeficiente es significativo en la regresión aumentada, $X_1$ y el error estructural están correlacionados (endogeneidad confirmada).}

\sol{53}{C}{\textbf{Tema:} Prueba de sobreidentificación de Sargan. \textbf{Explicación:} La hipótesis nula de la prueba de Sargan / Hansen es que todos los instrumentos son exógenos y válidos. Como el estadístico calculado $1.20 < 5.99$, no se rechaza $H_0$, respaldando la validez del instrumental utilizado.}

\sol{54}{A}{\textbf{Tema:} Trampa de las variables ficticias. \textbf{Explicación:} La suma de las 4 variables dummy es exactamente igual al vector columna de unos de la constante ($D_1 + D_2 + D_3 + D_4 = 1$), provocando dependencia lineal exacta. Para evitarlo se debe omitir una categoría base o suprimir la constante.}

\sol{55}{D}{\textbf{Tema:} Efecto marginal con variable de interacción. \textbf{Explicación:} La derivada parcial respecto a la educación es $\frac{\partial \ln(W)}{\partial Educ} = \beta_1 + \beta_3 Mujer$. Para hombres ($Mujer=0$) es $\beta_1$; para mujeres ($Mujer=1$) es $\beta_1 + \beta_3$.}

\sol{56}{C}{\textbf{Tema:} Propiedad de regresión estandarizada. \textbf{Explicación:} Al tener todas las variables media cero, la ecuación evaluada en las medias exige $\bar{Y} = \hat{\beta}_0 + \sum \hat{\beta}_j \bar{X}_j \implies 0 = \hat{\beta}_0 + 0 \implies \hat{\beta}_0 = 0$.}

\sol{57}{B}{\textbf{Tema:} Propiedades de la Hat Matrix. \textbf{Explicación:} $H$ proyecta el vector $Y$ sobre el subespacio generado por las columnas de $X$ ($\hat{Y} = HY$). Como operador de proyección ortogonal, es simétrica e idempotente, y su traza es igual al rango de $X$, que es $k$.}

\sol{58}{C}{\textbf{Tema:} Ortogonalidad de residuos MCO. \textbf{Explicación:} Las condiciones de primer orden de MCO exigen $\frac{\partial (e'e)}{\partial \beta} = -2X'e = 0 \implies X'e = 0$. Si hay intercepto, una columna de $X$ es de unos, garantizando $\sum e_i = 0$ y covarianza muestral nula con los regresores.}

\sol{59}{A}{\textbf{Tema:} Descomposición de varianza e intercepto. \textbf{Explicación:} El término cruzado $\sum (\hat{Y}_i - \bar{Y})e_i$ se anula únicamente si el modelo contiene constante (o si los residuos suman cero). Sin constante, $STC \neq SEC + SCR$ y el $R^2$ estándar carece de sentido matemático.}

\sol{60}{C}{\textbf{Tema:} Cálculo de intervalo de confianza. \textbf{Explicación:} $IC_{95\%} = \hat{\beta}_1 \pm t_{crit} \times EE(\hat{\beta}_1) = 2.50 \pm 2.00 \times 0.25 = 2.50 \pm 0.50 = [2.00, \; 3.00]$.}

\sol{61}{C}{\textbf{Tema:} Varianza del error de predicción individual. \textbf{Explicación:} La predicción individual incorpora dos fuentes de incertidumbre: la estimación de los parámetros $\beta$ (varianza $\sigma^2 x_0'(X'X)^{-1}x_0$) más la varianza de la nueva perturbación $u_0$ que afectará al individuo ($\\sigma^2$), sumando $\sigma^2 [1 + x_0'(X'X)^{-1}x_0]$.}

\sol{62}{B}{\textbf{Tema:} Comparación entre AIC y BIC. \textbf{Explicación:} La penalización del AIC es $2k$, mientras que la del BIC es $k \ln(n)$. Como para $n \ge 8$ se cumple $\ln(n) > 2$, el BIC castiga más drásticamente la sobreparametrización y es un criterio consistente para identificar el modelo verdadero.}

\sol{63}{D}{\textbf{Tema:} Prueba de normalidad de Jarque-Bera. \textbf{Explicación:} Bajo la hipótesis nula de normalidad, $S \approx 0$ y $K \approx 3$. El estadístico de Jarque-Bera combina las desviaciones cuadráticas de asimetría y exceso de curtosis, siguiendo asintóticamente una distribución $\chi^2$ con 2 grados de libertad.}

\sol{64}{B}{\textbf{Tema:} Prueba RESET de Ramsey. \textbf{Explicación:} Si la forma funcional lineal es correcta, las potencias de los valores predichos ($\hat{Y}^2, \hat{Y}^3$) no deberían tener poder explicativo residual. Un estadístico $F$ significativo indica no linealidad o sesgo por variables omitidas.}

\sol{65}{D}{\textbf{Tema:} Propiedad de MCO restringido. \textbf{Explicación:} Imponer restricciones reduce el espacio de búsqueda en la minimización de residuos. Por ende, la suma de residuos al cuadrado restringida nunca puede ser menor que la no restringida: $SCR_R \ge SCR_{NR}$ (igualdad solo si la muestra satisface exactamente la restricción).}

\subsection*{Bloque C. Casos con implicancias regulatorias y sector saneamiento}

\sol{66}{B}{\textbf{Tema:} Modelación de demanda de agua en SUNASS. \textbf{Explicación:} La forma funcional log-log permite interpretar inmediatamente los coeficientes estimados como las elasticidades directas de consumo frente a cambios tarifarios y variaciones en el ingreso de los hogares, parámetros esenciales para diseñar la estructura tarifaria quinquenal.}

\sol{67}{B}{\textbf{Tema:} Endogeneidad en tarifas por bloques. \textbf{Explicación:} Bajo esquemas de bloques crecientes, la tarifa marginal no es exógena: depende del volumen consumido. Si un choque no observado eleva el consumo ($u > 0$), el usuario salta de bloque tarifario ($P$ sube), generando correlación positiva espuria entre precio y error.}

\sol{68}{A}{\textbf{Tema:} Instrumentación en demanda de agua. \textbf{Explicación:} La tarifa del primer bloque básico es fijada exógenamente por el regulador y es independiente del consumo discrecional final del hogar, funcionando como un instrumento válido que correlaciona con la tarifa promedio sin compartir choques de demanda.}

\sol{69}{B}{\textbf{Tema:} Frontera estocástica en regulación de EPS. \textbf{Explicación:} SFA especifica el residuo como $\varepsilon_i = v_i + u_i$, donde $v_i \sim N(0, \sigma_v^2)$ es ruido estocástico blanco (clima, geografía) y $u_i \ge 0$ es ineficiencia técnica pura. Esto evita atribuir a mala gestión costos derivados de choques exógenos adversos.}

\sol{70}{B}{\textbf{Tema:} Economías de densidad en saneamiento. \textbf{Explicación:} Las economías de densidad de clientes/producto miden cómo responde el costo al expandir el volumen suministrado dentro de una red de distribución ya existente (sin extender kilómetros de tubería), siendo típicamente crecientes en servicios de saneamiento.}

\sol{71}{D}{\textbf{Tema:} Heterocedasticidad por tamaño de empresa. \textbf{Explicación:} En muestras con disparidad extrema de tamaño (Sedapal atiende a más de 10 millones de habitantes mientras que pequeñas EPS atienden a 20 mil), la variabilidad de los errores es función directa del tamaño de la empresa, requiriendo estimación ponderada o errores robustos.}

\sol{72}{D}{\textbf{Tema:} Errores robustos HC3 en muestras pequeñas. \textbf{Explicación:} Davidson y MacKinnon demostraron que en muestras con menos de 100 observaciones, el estimador HC0 de White subestima la verdadera varianza en puntos de alto leverage. El estimador HC3 divide los residuos al cuadrado entre $(1 - h_{ii})^2$, corrigiendo el sesgo muestral.}

\sol{73}{A}{\textbf{Tema:} Modelación SARIMA en hidrología urbana. \textbf{Explicación:} Las series hidrológicas presentan alta memoria temporal y patrones estacionales marcados (estiaje y avenidas). Los modelos SARIMA $(p,d,q)(P,D,Q)_{12}$ capturan la dinámica de corto plazo y la estacionalidad anual del recurso hídrico.}

\sol{74}{B}{\textbf{Tema:} Newey-West en series de tiempo climáticas. \textbf{Explicación:} Los choques climáticos como El Niño generan correlaciones residuales que persisten a lo largo de múltiples trimestres. Los errores HAC de Newey-West blindan la inferencia estadística evitando rechazos espurios de significancia causal.}

\sol{75}{B}{\textbf{Tema:} Sesgo de atrición en paneles de medidores. \textbf{Explicación:} Si la salida de medidores del panel no es un proceso puramente aleatorio (ej. se rompen más rápido en casas de consumo masivo o con mayor presión de agua), la muestra superviviente se vuelve no representativa, sesgando los coeficientes de demanda.}

\sol{76}{C}{\textbf{Tema:} Diferencias en Diferencias en saneamiento. \textbf{Explicación:} Diff-in-Diff controla por factores macroeconómicos y climáticos temporales comunes que afectan a toda la ciudad: calcula la diferencia antes-después en el grupo tratado y le resta la diferencia antes-después del grupo de control: $\hat{\delta}_{DD} = (\bar{Y}_{T,post} - \bar{Y}_{T,pre}) - (\bar{Y}_{C,post} - \bar{Y}_{C,pre})$.}

\sol{77}{D}{\textbf{Tema:} Supuesto de tendencias paralelas. \textbf{Explicación:} Si en el periodo previo al tratamiento el grupo de control y tratamiento ya exhibían trayectorias divergentes, la diferencia observada post-tratamiento no puede atribuirse con rigor causal a la instalación del medidor, invalidando la evaluación.}

\sol{78}{A}{\textbf{Tema:} Regresión Discontinua (RDD) en SISFOH. \textbf{Explicación:} Alrededor del punto de corte $c_0$, los hogares justo por encima y justo por debajo del umbral tienen características socioeconómicas prácticamente idénticas (asignación cuasialeatoria), permitiendo estimar el efecto causal del subsidio tarifario de forma muy robusta.}

\sol{79}{B}{\textbf{Tema:} Limitaciones del Modelo Lineal de Probabilidad. \textbf{Explicación:} El MLP ($P(Y=1|X) = X\beta$) genera predicciones absurdas ($P < 0$ o $P > 1$) en los extremos de la distribución y su error tiene varianza heterocedástica $Var(u_i|X) = X\beta(1 - X\beta)$. Logit y Probit confinan la probabilidad entre 0 y 1 mediante funciones de distribución acumuladas.}

\sol{80}{D}{\textbf{Tema:} Multicolinealidad en infraestructura de EPS. \textbf{Explicación:} Las EPS con más kilómetros de red tienen inevitablemente más conexiones domiciliarias y producen más volumen de agua. Esta colinealidad infla las varianzas de los coeficientes individuales, aunque el modelo conjunto prediga con gran exactitud el costo total.}

\sol{81}{A}{\textbf{Tema:} Causalidad inversa en ANF y costos. \textbf{Explicación:} Existe retroalimentación simultánea entre el gasto en mantenimiento de redes y las fugas físicas de agua: más fugas demandan más reparaciones, pero menos mantenimiento provoca más roturas de tuberías, requiriendo variables instrumentales para identificar el efecto causal.}

\sol{82}{D}{\textbf{Tema:} Corrección de selección de Heckman. \textbf{Explicación:} Los hogares con conexión domiciliaria no son una muestra aleatoria de la población urbana (autoselección). Estimar MCO ignorando a los no conectados introduce sesgo de selección. El método de Heckman estima primero la probabilidad de tener conexión (Probit) e inserta la Razón Inversa de Mills en la ecuación de gasto.}

\sol{83}{C}{\textbf{Tema:} Ventaja de Efectos Fijos en datos de panel de EPS. \textbf{Explicación:} Cada EPS opera bajo condiciones geográficas y topográficas únicas que no varían con el tiempo. El estimador Within de Efectos Fijos elimina $\alpha_i$ diferenciando respecto a la media individual, evitando sesgos por variables omitidas territoriales.}

\sol{84}{A}{\textbf{Tema:} Prueba de Hausman en datos de panel. \textbf{Explicación:} Bajo la hipótesis nula, ambos estimadores son consistentes pero RE es más eficiente. Si $H_0$ se rechaza, la heterogeneidad no observada $\alpha_i$ está correlacionada con los regresores, haciendo que RE sea sesgado e inconsistente, por lo que la teoría exige adoptar Efectos Fijos.}

\sol{85}{D}{\textbf{Tema:} Estimación de la Beta CAPM en WACC regulatorio. \textbf{Explicación:} La Beta mide el riesgo sistemático no diversificable del negocio de saneamiento. Se estima regresionando los retornos bursátiles de empresas del sector sobre el índice de mercado (S\&P 500) y ajustando por apalancamiento financiero.}

\sol{86}{D}{\textbf{Tema:} Cálculo de Beta de Hamada en regulación. \textbf{Explicación:} $\beta_L = 0.50 \times [1 + (1 - 0.30) \times 0.60] = 0.50 \times [1 + 0.70 \times 0.60] = 0.50 \times [1 + 0.42] = 0.50 \times 1.42 = 0.71$.}

\sol{87}{C}{\textbf{Tema:} Modelos de conteo en roturas de tuberías. \textbf{Explicación:} Las fallas de red son eventos discretos no negativos. El modelo de Poisson modela conteos pero asume equi-dispersión ($Var = Media$). Si los datos presentan sobredispersión por incidentes en racimo, el modelo Binomial Negativo introduce un parámetro de dispersión gamma para inferencia correcta.}

\sol{88}{B}{\textbf{Tema:} Detección econométrica de cárteles en licitaciones. \textbf{Explicación:} Los cárteles coordinan ofertas de resguardo que se agrupan en torno a márgenes prefijados, generando una distribución de posturas con dispersión artificialmente baja y patrones de simetría incompatibles con la competencia no cooperativa.}

\sol{89}{D}{\textbf{Tema:} Econometría de Nowcasting en saneamiento. \textbf{Explicación:} El Nowcasting (predicción del presente inmediato) utiliza modelos de factores dinámicos o MIDAS para explotar indicadores de alta frecuencia disponibles de inmediato, permitiendo anticipar caídas en la cobranza de las EPS antes de que se publiquen los estados financieros oficiales.}

\sol{90}{D}{\textbf{Tema:} Poder estadístico y errores Tipo I y II en calidad de agua. \textbf{Explicación:} El diseño del muestreo bacteriológico es una decisión econométrica/estadística crucial: rechazar erróneamente el cumplimiento impone costos judiciales injustos a la EPS (Tipo I), mientras que no detectar contaminación expone a la población a epidemias diarreicas (Tipo II), requiriendo un poder estadístico ($1 - \beta$) riguroso.}

\subsection*{Bloque D. Ampliación: Teoría y Aplicaciones Prácticas (91 a 180)}

\sol{91}{A}{La matriz $M = I_n - P_X$ proyecta sobre el complemento ortogonal de $X$: es simétrica ($M'=M$), idempotente ($M^2=M$), aniquila a $X$ ($MX=0$) y $\text{tr}(M)=n-k$.}
\sol{92}{B}{El teorema FWL demuestra que $\hat{\beta}_1 = (X_1' M_2 X_1)^{-1} X_1' M_2 Y$, es decir, MCO sobre las variables previamente depuradas de la influencia lineal de $X_2$.}
\sol{93}{C}{Sin intercepto en $X$, la condición de primer orden $X'\hat{u}=0$ no garantiza $\sum \hat{u}_i = 0$, por lo que $SCT \neq SCE + SCR$ y $1 - SCR/SCT$ puede ser negativo.}
\sol{94}{D}{Por la ley de expectativas iteradas, $E[u|X]=0 \Rightarrow E[X'u]=0$, pero no al revés. $E[u|X]=0$ da insesgadez exacta; $E[X'u]=0$ garantiza consistencia asintótica.}
\sol{95}{B}{Bajo normalidad, MCO coincide con MV y alcanza la cota de Cramér-Rao, siendo el Mejor Estimador Insesgado (MIE/BUE) incluso frente a estimadores no lineales.}
\sol{96}{C}{Si $X_2^* = c X_2$, entonces $\hat{\beta}_2^* = \hat{\beta}_2/c$ y $ee(\hat{\beta}_2^*) = ee(\hat{\beta}_2)/c$; por tanto, $t_2$, $\hat{Y}$ y $R^2$ permanecen exactamente invariantes.}
\sol{97}{D}{$u_i = Y_i - X_i'\beta$ es inobservable y esférico ($\sigma^2 I$), en tanto $\hat{u} = Mu$ es observable y tiene matriz de varianzas $\sigma^2 M$ (ni independiente ni homocedástica).}
\sol{98}{A}{Puede demostrarse algebraicamente que añadir un regresor reduce la varianza residual estimada $s^2 = SCR/(n-k)$ (y por ende sube $\bar{R}^2$) si y solo si $|t_{k+1}| > 1$.}
\sol{99}{C}{Al pasar de $D=0$ a $D=1$, $\frac{Y_1 - Y_0}{Y_0} = e^{\delta} - 1$. Así, el cambio porcentual exacto es $100(e^{\delta}-1)\%$; la regla $100\delta\%$ es solo una aproximación de primer orden.}
\sol{100}{D}{Derivando respecto a $X$: $\frac{\partial E[Y|X]}{\partial X} = \beta_1 + 2\beta_2 X$. Igualando a cero se obtiene el punto crítico $X^* = -\beta_1/(2\beta_2)$, que es máximo pues $2\beta_2 < 0$.}
\sol{101}{A}{El cuadrado de una variable $t_{n-k}$ sigue exactamente una distribución $F_{1,n-k}$. Por ello, para $q=1$ restricción lineal, el contraste $F$ y el contraste $t$ bilateral son idénticos.}
\sol{102}{B}{Como $ee(\hat{\beta}_j) \propto 1/\sqrt{n}$, en muestras muy grandes cualquier efecto real distinto de cero —por diminuto que sea— genera $|t|$ muy grande y $p$-valor cercano a $0$.}
\sol{103}{D}{La prueba de Chow contrasta si tanto el intercepto como todas las pendientes permanecen constantes entre submuestras, lo cual equivale a probar conjuntamente la dummy de periodo y todas sus interacciones.}
\sol{104}{A}{En el modelo lineal normal con restricciones lineales se cumple exactamente $W \ge LR \ge LM$, y asintóticamente los tres estadísticos son equivalentes bajo $H_0$ con distribución $\chi^2_q$.}
\sol{105}{B}{Sobre-especificar el modelo añadiendo una variable con $\beta=0$ no introduce sesgo ni inconsistencia en MCO, pero reduce eficiencia (infla las varianzas) si correlaciona con los demás regresores.}
\sol{106}{C}{Con $\text{Cov}(X_1, X_2)=0$, el sesgo de variable omitida en $\hat{\beta}_1$ es cero; sin embargo, $\beta_2 X_2$ pasa al término de error, inflando $s^2$ y sesgando al alza los errores estándar.}
\sol{107}{A}{Ambos criterios se minimizan. Como para $n \ge 8$ se cumple $\ln(n) > 2$, BIC penaliza más la complejidad que AIC, favorece modelos más parsimoniosos y es consistente en selección.}
\sol{108}{B}{Con $n$ fijo, reducir $\alpha$ estrecha la región de rechazo, dificultando rechazar $H_0$: baja el Error Tipo I, pero sube el Error Tipo II ($\beta$) y disminuye la potencia ($1-\beta$).}
\sol{109}{C}{La multicolinealidad infla las varianzas individuales (bajando los $|t|$), pero no afecta la capacidad explicativa conjunta ($R^2$ y $F$ global altos) ni viola el carácter MELI de MCO.}
\sol{110}{D}{$R_j^2$ es el $R^2$ de regresar $X_j$ contra las demás explicativas; $1/(1-R_j^2)$ es el $VIF_j$, que mide cuánto se infla la varianza de $\hat{\beta}_j$ por colinealidad.}
\sol{111}{B}{La prueba de White detecta divergencias entre la matriz de varianzas clásica y la robusta, lo cual ocurre tanto ante heterocedasticidad pura como ante mala especificación funcional.}
\sol{112}{C}{Goldfeld-Quandt ordena según $X_j$, elimina $c$ observaciones intermedias para acentuar la diferencia de varianzas, estima MCO en ambos extremos y contrasta $F = SCR_2/SCR_1$.}
\sol{113}{D}{Al sustituir $\Omega$ por una estimación $\hat{\Omega}$, FGLS pierde la garantía exacta de insesgadez en muestras finitas, pero conserva consistencia y eficiencia asintótica sobre MCO.}
\sol{114}{A}{Los procesos ARCH/GARCH modelan agrupamientos de volatilidad (*volatility clustering*): la varianza condicional $\sigma_t^2$ varía en el tiempo aunque la varianza incondicional sea finita y constante.}
\sol{115}{C}{Como $DW \approx 2(1 - \hat{\rho})$, cuando $\hat{\rho} < 0$ se tiene $DW > 2$, acercándose a $4$ cuando $\hat{\rho} \to -1$.}
\sol{116}{D}{Con $Y_{t-1}$ en el lado derecho, DW se sesga hacia 2; se emplea el estadístico $h$ de Durbin (para $AR(1)$) o preferentemente Breusch-Godfrey (válido para cualquier orden $p$).}
\sol{117}{A}{Cochrane-Orcutt descarta $t=1$ al cuasi-diferenciar ($y_t - \hat{\rho}y_{t-1}$). Prais-Winsten preserva $t=1$ escalando por $\sqrt{1-\hat{\rho}^2}$, lo que mejora la eficiencia en muestras pequeñas.}
\sol{118}{B}{Al igual que los errores robustos de White, Newey-West no modifica $\hat{\beta}_{MCO}$; solo reemplaza la matriz de varianzas por un estimador consistente ante heterocedasticidad y autocorrelación.}
\sol{119}{D}{Un error clásico en $Y$ se absorbe en $u$, manteniendo insesgadez con mayor varianza. Un error clásico en $X$ correlaciona el regresor observado con el error compuesto, sesgando $\hat{\beta}_1$ hacia cero.}
\sol{120}{A}{La condición de orden (contar exclusiones: $K - k \ge m - 1$) es necesaria mas no suficiente; la condición de rango sobre la matriz de parámetros reducidos es necesaria y suficiente.}
\sol{121}{B}{Imbens y Angrist (1994) demostraron que con efectos heterogéneos, validez del instrumento y monotonicidad, el estimador VI identifica el LATE: el efecto causal promedio para los *compliers*.}
\sol{122}{C}{Con instrumentos débiles ($\text{Cov}(Z,X) \approx 0$), MC2E tiene una distribución no normal, varianza enorme y un sesgo en muestras finitas que lo arrastra hacia el estimador sesgado de MCO.}
\sol{123}{A}{Con $L = m$ (exactamente identificado), los residuos VI cumplen $Z'\hat{u}_{VI} = 0$ por definición algebraica ($0$ grados de sobreidentificación), por lo que la exclusión debe defenderse teóricamente.}
\sol{124}{B}{MC2E es un caso particular de GMM que usa la matriz de pesos óptima solo bajo homocedasticidad. Con heterocedasticidad y sobreidentificación, GMM en dos etapas es asintóticamente más eficiente.}
\sol{125}{C}{La prueba de Sargan asume errores homocedásticos; bajo heterocedasticidad o errores agrupados, el estadístico $J$ de Hansen (basado en GMM) mantiene distribución asintótica $\chi^2_{L-m}$.}
\sol{126}{D}{Controlar por variables determinadas después del tratamiento (*bad controls*) introduce sesgo de selección al condicionar en un resultado o colisionador, distorsionando el efecto causal total.}
\sol{127}{B}{La transformación intragrupo resta la media temporal de cada individuo. Si $z_{it} = z_i$ para todo $t$, entonces $z_{it} - \bar{z}_i = 0$ y queda colinealmente absorbida por el efecto fijo $\alpha_i$.}
\sol{128}{C}{Por el Teorema de Frisch-Waugh-Lovell, proyectar fuera las dummies individuales produce exactamente las desviaciones respecto a la media individual; por ello $\hat{\beta}_{within} \equiv \hat{\beta}_{LSDV}$.}
\sol{129}{D}{Si $u_{it}$ es ruido blanco, FE es más eficiente que FD (pues $\Delta u_{it}$ tiene autocorrelación $-0.5$). Si $u_{it}$ es caminata aleatoria, $\Delta u_{it} = e_{it}$ es serialmente no correlacionado y FD es más eficiente.}
\sol{130}{A}{Como $\theta = 1 - \sqrt{\frac{\sigma_u^2}{\sigma_u^2 + T\sigma_{\alpha}^2}}$, si $T\sigma_{\alpha}^2 \gg \sigma_u^2$ entonces $\theta \to 1$ y RE se aproxima a FE; si $\sigma_{\alpha}^2 \to 0$, $\theta \to 0$ y RE tiende a MCO agrupado.}
\sol{131}{C}{Agrupar (*cluster*) los errores estándar a nivel de unidad (EPS) garantiza inferencia asintóticamente válida ante heterocedasticidad y correlación serial arbitraria dentro de cada clúster.}
\sol{132}{D}{En paneles cortos con $y_{it-1}$, la media $\bar{u}_i$ contiene $u_{it-1}$, correlacionando $\tilde{y}_{it-1}$ con $\tilde{u}_{it}$ (sesgo de Nickell). Arellano-Bond usa primeras diferencias e instrumentos en niveles rezagados.}
\sol{133}{A}{FE es robusto a la atrición si la selección depende de regresores $X_{it}$ o de $\alpha_i$ (ya que $\alpha_i$ se elimina), pero se sesga si la salida depende de los choques idiosincráticos $u_{it}$.}
\sol{134}{B}{En un proceso TS ($y_t = \alpha + \delta t + u_t$ con $u_t$ estacionario), los shocks mueren y $y_t$ vuelve a la tendencia. En un proceso DS ($y_t = \delta + y_{t-1} + \varepsilon_t$), cada shock se acumula permanentemente.}
\sol{135}{D}{Un proceso $AR(p)$ es estacionario si todas las raíces de $\phi(z)=1-\phi_1 z - \dots - \phi_p z^p = 0$ satisfacen $|z| > 1$ (en el caso $AR(1)$, $1-\phi_1 z=0 \Rightarrow |z|=1/|\phi_1|>1 \Leftrightarrow |\phi_1|<1$).}
\sol{136}{A}{En un proceso $MA(q)$, la memoria se corta en el rezago $q$ (FAC $= 0$ para $k > q$) y la FACP decae gradualmente; en un $AR(p)$ ocurre exactamente el patrón simétrico inverso.}
\sol{137}{B}{La prueba DF simple asume errores no correlacionados. La versión Aumentada (ADF) añade rezagos de $\Delta y_t$ para limpiar la autocorrelación residual, manteniendo valores críticos de Dickey-Fuller.}
\sol{138}{C}{El Teorema de Representación de Granger establece que dos series $I(1)$ están cointegradas si y solo si admiten un MCE, implicando además causalidad de Granger en al menos una dirección.}
\sol{139}{A}{Si $y_{t-1}$ excede su valor de equilibrio ($u_{t-1} > 0$), la corrección requiere que $\Delta y_t$ baje; por ello se exige $\alpha < 0$ (típicamente $-1 < \alpha < 0$ para convergencia monótona).}
\sol{140}{B}{La causalidad de Granger evalúa precedencia predictiva (contrastada mediante prueba $F$/Wald sobre los rezagos de $X$ en la ecuación de $Y$), no causalidad estructural filosófica.}
\sol{141}{C}{Cholesky impone una estructura recursiva triangular inferior sobre los impactos contemporáneos; si la matriz de covarianzas residuales no es diagonal, cambiar el orden de las variables altera las IRF.}
\sol{142}{D}{Como $\frac{\partial P(Y=1|X)}{\partial X_j} = g(X'\beta)\beta_j$ y $g(\cdot)>0$, el signo coincide con $\beta_j$, pero la magnitud varía según $X$; se resume mediante el efecto en la media (MEM) o el promedio (AME).}
\sol{143}{B}{En Logit, $\ln\left(\frac{p}{1-p}\right) = X'\beta$, de modo que aumentar $X_j$ en una unidad multiplica los *odds* $\frac{p}{1-p}$ exactamente por $e^{\beta_j}$.}
\sol{144}{C}{Con censura (solución de esquina o límite de medición) conocemos $X_i$ de toda la muestra aunque $Y_i$ se apile en el límite; con truncamiento se excluyen de la base los individuos fuera del rango.}
\sol{145}{D}{Poisson impone equidispersión ($\text{Var}=E$). Bajo sobredispersión, Poisson cuasi-MV sigue siendo consistente para $\beta$ si la media está bien especificada, pero los $ee$ estándar son demasiado pequeños; Binomial Negativa modela explícitamente la sobredispersión.}
\sol{146}{A}{Como en modelos discretos $0 > \ln L_{ur} \ge \ln L_0$, el cociente $\ln L_{ur}/\ln L_0 \in (0, 1]$ y $R^2_{McF} = 1 - \ln L_{ur}/\ln L_0 \in [0, 1)$, evaluando mejora relativa en log-verosimilitud.}
\sol{147}{C}{PSM asume que, condicionando en $X$ (o en $p(X)$), el tratamiento es independiente de los resultados potenciales (CIA) y que existe traslape de probabilidades $0 < P(D=1|X) < 1$ (soporte común).}
\sol{148}{D}{En *Fuzzy RDD*, el indicador de cruzar el umbral actúa como instrumento válido para el tratamiento efectivo cerca del corte: $\hat{\tau}_{FRD} = \frac{\lim_{x\downarrow c}E[Y|x] - \lim_{x\uparrow c}E[Y|x]}{\lim_{x\downarrow c}E[D|x] - \lim_{x\uparrow c}E[D|x]}$.}
\sol{149}{A}{La validez de RDD exige que los individuos no puedan manipular con precisión su puntaje alrededor del corte; McCrary contrasta la continuidad de la densidad de la *running variable* en el umbral.}
\sol{150}{B}{Goodman-Bacon (2021) y Callaway-Sant'Anna muestran que con adopción escalonada y efectos dinámicos heterogéneos, TWFE realiza comparaciones 'prohibidas' usando ya tratados como controles, pudiendo asignar pesos negativos.}
\sol{151}{D}{DEA construye una frontera determinística por programación lineal (sensible a outliers/ruido); SFA estima una frontera paramétrica descomponiendo la desviación en ruido simétrico y término de ineficiencia.}
\sol{152}{A}{En SFA de costos (minimización), $v_i \sim N(0, \sigma_v^2)$ es ruido aleatorio bidireccional y $u_i \ge 0$ (seminormal, exponencial o normal truncada) mide el sobrecosto por ineficiencia.}
\sol{153}{B}{En Heckman en dos etapas, $\hat{\lambda}_i = \phi(z_i'\hat{\gamma})/\Phi(z_i'\hat{\gamma})$ corrige la media condicional truncada; si su coeficiente es significativamente distinto de cero, existe sesgo de selección.}
\sol{154}{C}{La recta MCO pasa por $(\bar{X}, \bar{Y})$: $\hat{\beta}_0 = \bar{Y} - \hat{\beta}_1\bar{X} = 85 - 2.5(20) = 35.0$. Luego, $\hat{Y} = 35.0 + 2.5(24) = 95.0$.}
\sol{155}{A}{$R^2 = 1 - \frac{100}{500} = 0.80$. Además, $\bar{R}^2 = 1 - \frac{SCR/(n-k)}{SCT/(n-1)} = 1 - \frac{100/40}{500/44} = 1 - 0.22 = 0.78$.}
\sol{156}{B}{Aplicando $\bar{R}^2 = 1 - (1 - R^2)\frac{n-1}{n-k} = 1 - (0.25)\frac{50}{45} = 1 - 0.2778 = 0.7222$.}
\sol{157}{C}{El estadístico es $t = \frac{\hat{\beta}_1 - \beta_{1,0}}{ee(\hat{\beta}_1)} = \frac{-0.64 - (-1.00)}{0.12} = \frac{0.36}{0.12} = 3.00$.}
\sol{158}{D}{Usando $F = \frac{(R^2_{NR} - R^2_R)/q}{(1 - R^2_{NR})/(n - k)} = \frac{(0.70 - 0.58)/3}{(1 - 0.70)/(66 - 6)} = \frac{0.04}{0.005} = 8.00$.}
\sol{159}{B}{Para $q = 1$ se cumple $F = t^2$, luego $|t| = \sqrt{16} = 4.0$. Así, $ee(\hat{\beta}_j) = \frac{|\hat{\beta}_j|}{|t|} = \frac{1.20}{4.0} = 0.300$.}
\sol{160}{C}{El cambio porcentual exacto es $100(e^{\hat{\delta}} - 1)\% = 100(e^{-0.20} - 1)\% = 100(0.8187 - 1)\% = -18.13\%$.}
\sol{161}{D}{Derivando respecto a $X$: $\frac{d\hat{Y}}{dX} = 12 - 0.60X = 0 \Rightarrow X^* = \frac{12}{0.60} = 20$ años (máximo porque $-0.60 < 0$).}
\sol{162}{A}{El efecto marginal es $\frac{\partial \hat{Y}}{\partial X} = 0.80 - 0.25D$. Evaluando en $D = 1$ se obtiene $0.80 - 0.25 = 0.55$ unidades.}
\sol{163}{C}{Como $VIF_3 = \frac{1}{1 - R_3^2} = 8$, se despeja $1 - R_3^2 = \frac{1}{8} = 0.125 \Rightarrow R_3^2 = 0.875$.}
\sol{164}{D}{El estadístico de Breusch-Godfrey es $LM = n \cdot R^2_{aux} = 100 \times 0.085 = 8.50$, el cual sigue asintóticamente una distribución $\chi^2_p = \chi^2_2$ bajo $H_0$.}
\sol{165}{A}{Despejando de $DW \approx 2(1 - \hat{\rho})$: $\hat{\rho} \approx 1 - \frac{DW}{2} = 1 - \frac{1.30}{2} = 1 - 0.65 = 0.35$.}
\sol{166}{B}{Bajo error de medición clásico en $X$, $\text{plim}\,\hat{\beta}_1 = \beta_1 \frac{\sigma_{X^*}^2}{\sigma_{X^*}^2 + \sigma_e^2} = 2.0 \times \frac{16}{16 + 4} = 2.0 \times 0.80 = 1.60$.}
\sol{167}{D}{El estimador de Wald es la razón entre el efecto intención de tratar y la primera etapa: $\hat{\beta}_{Wald} = \frac{\bar{Y}_1 - \bar{Y}_0}{\bar{X}_1 - \bar{X}_0} = \frac{42 - 36}{0.75 - 0.25} = \frac{6}{0.50} = 12.0$.}
\sol{168}{A}{El sesgo de variable omitida es $\text{Sesgo}(\hat{\beta}_1) = \beta_2 \frac{\text{Cov}(X_1, X_2)}{\text{Var}(X_1)} = 4.0 \times \frac{-3.0}{10.0} = -1.20$.}
\sol{169}{B}{$\hat{\delta}_{DiD} = (\bar{Y}_{T,1} - \bar{Y}_{T,0}) - (\bar{Y}_{C,1} - \bar{Y}_{C,0}) = (17.5 - 22.0) - (18.5 - 20.0) = -4.5 - (-1.5) = -3.0\ m^3$.}
\sol{170}{C}{En el modelo Logit, $\hat{P}(Y=1|X) = \frac{1}{1 + e^{-X'\hat{\beta}}} = \frac{1}{1 + e^{-1}} = \frac{1}{1.3679} \approx 0.7311$.}
\sol{171}{A}{El *odds ratio* es $e^{\hat{\beta}_1} = e^{0.40} \approx 1.4918$, lo que equivale a un incremento porcentual en los *odds* de $(1.4918 - 1) \times 100\% = 49.18\%$.}
\sol{172}{B}{El pseudo-$R^2$ de McFadden es $R^2_{McF} = 1 - \frac{\ln L_{ur}}{\ln L_0} = 1 - \frac{-120}{-200} = 1 - 0.60 = 0.40$.}
\sol{173}{C}{Para un $AR(1)$ estacionario: $E[y_t] = \frac{c}{1-\phi} = \frac{6}{0.4} = 15$, y $\text{Var}(y_t) = \frac{\sigma_\varepsilon^2}{1-\phi^2} = \frac{16}{1-0.36} = \frac{16}{0.64} = 25$.}
\sol{174}{D}{En un proceso $AR(1)$ estacionario, la autocorrelación de orden $k$ cumple $\rho_k = \phi^k$. Para $k = 3$: $\rho_3 = (0.8)^3 = 0.512$.}
\sol{175}{B}{Para un $MA(1)$: $\text{Var}(y_t) = \sigma_\varepsilon^2(1 + \theta^2) = 4(1 + 0.25) = 5.0$ y $\rho_1 = \frac{\theta}{1 + \theta^2} = \frac{0.5}{1.25} = 0.40$.}
\sol{176}{C}{Con $\Delta x_t = 0$ y desequilibrio $(y_{t-1} - y^*_{t-1}) = +8$, el ajuste esperado es $E[\Delta y_t] = -0.25(8) = -2.00$ unidades (corrige $25\%$ del desvío por periodo).}
\sol{177}{D}{$AIC = 2.00 + \frac{2(5)}{100} = 2.100$. $BIC = 2.00 + \frac{5(4.605)}{100} = 2.00 + 0.230 = 2.230$.}
\sol{178}{A}{Sustituyendo en la fórmula: $JB = \frac{120}{6}\left[(0.5)^2 + \frac{(4.0 - 3.0)^2}{4}\right] = 20(0.25 + 0.25) = 10.00$.}
\sol{179}{C}{En un modelo lin-log, $\Delta Q \approx \left(\frac{\hat{\beta}_1}{100}\right) \%\Delta X = \left(\frac{30}{100}\right)(10) = 3.00\ m^3$.}
\sol{180}{D}{La media temporal de la EPS es $\bar{x}_i = \frac{12 + 18 + 21 + 25}{4} = 19.0$. Por tanto, para $t = 3$ se obtiene $\tilde{x}_{i3} = 21 - 19 = +2.0$.}
